\section{Discussion}

The performance of the PCA foreground removal method was evaluated by comparing the recovered residual signal to the HI signal. By removing $N_{FG} = 3$, the majority of the foreground, which is four ordered of magnitude larger than the 21\,cm signal, was removed. 

The ratio plot in \Cref{fig:cylindrical_power_spectrum_hi_signal} seen on the right reveals the outcome of the foreground cleaning. Due to the foregrounds being spectrally smooth, they occupy the low-$k_\parallel$ region of the power spectrum which corresponds to the large-scale frequency of the signal. However, it can be seen that low values of $k_\parallel$ have significant signal loss observed below $k_\parallel = 0.1$, indicating the 21\,cm signal cannot be separated from the foreground contribution in these modes. %Similarly, there is overlap at low $k_\perp$ due to the smoothness of the foregrounds. 

To address this, minimum thresholds of $k_\perp = 0.025$ and $k_\parallel = 0.1$ were defined and foreground cleaning of the signal was redone excluding these values to ensure minimal 21\,cm signal loss. After applying these thresholds, the resulting 1D power spectrum can be seen in \Cref{fig:log_power_spectrum_hi_signal}. The resulting residual signal shows strong agreement with the 21\,cm signal at the high $k$. This suggests that PCA is effective at recovering the high-frequency fluctuations of the 21\,cm signal and is physically expected since high $k_\parallel$ correspond to small-scale frequency structure that the foregrounds do not contaminate. However, the residual signal is slightly above the expected HI signal at lower $k_\parallel$, suggesting there are residual foregrounds that PCA did not fully capture and the cleaning is incomplete.

%Therefore, while PCA is an effective blind foreground removal method, its main limitation is the removal of entire modes which results in the observed signal loss at low $k_\parallel$. 

%Therefore, while PCA is an effective blind foreground removal method for recovering high-frequency 21\,cm fluctuations, its main limitation is the loss of entire modes at low k∥ and incomplete foreground separation near the foreground-dominated boundary.